\nopagebreak

La región está limitada por la curva \(y = 5 - |x|\), el eje \(x\) y las rectas
\(x=-5\) y \(x=5\).  

Debido a la simetría, podemos integrar de \(0\) a \(5\) y multiplicar por 2:
\[
A = 2 \int_{0}^{5} (5 - x) \, dx.
\]

Dividimos el intervalo \([0,5]\) en \(n\) partes iguales:
\[
\Delta x = \frac{5-0}{n} = \frac{5}{n}.
\]

Tomando puntos derechos:
\[
x_i = 0 + \frac{5i}{n} = \frac{5i}{n}.
\]

El área se expresa como límite de sumas de Riemann:
\[
A = 2 \lim_{n \to \infty} \sum_{i=1}^{n} \left(5 - x_i\right) \Delta x.
\]

Sustituyendo \(x_i\):
\[
A = 2 \lim_{n \to \infty} \sum_{i=1}^{n} \left(5 - \frac{5i}{n}\right) \frac{5}{n}.
\]

Simplificando:
\[
A = 2 \lim_{n \to \infty} \sum_{i=1}^{n} \left(\frac{25}{n} - \frac{25i}{n^2}\right).
\]

Separando sumas:
\[
A = 2 \lim_{n \to \infty} \left[
\frac{25}{n} \sum_{i=1}^{n} 1 - \frac{25}{n^2} \sum_{i=1}^{n} i
\right].
\]

Usando:
\[
\sum_{i=1}^{n} 1 = n, \quad \sum_{i=1}^{n} i = \frac{n(n+1)}{2},
\]

se obtiene:
\[
A = 2 \lim_{n \to \infty} \left[ 25 - \frac{25(n+1)}{2n} \right]
= 2 \cdot \frac{25}{2} = 25.
\]

Finalmente:
\[
\boxed{A = 25}
\]
